\SetPageTheme{story}
\PageHeader{Lectia 5}{Structura repetitiva \texttt{cat timp}}{L5 • P1}

\begin{missionbox}[De ce ne trebuie repetitie]
  Pana acum am invatat sa extragem o cifra. Dar un numar are mai multe cifre. Daca vrem sa le prelucram pe toate, nu putem scrie aceeasi instructiune de zece ori sau de o suta de ori.

  Avem nevoie de un mecanism care sa repete aceeasi actiune pana cand numarul nu mai are nimic de oferit. Acesta este rolul structurii \texttt{cat timp}.
\end{missionbox}

\begin{conceptbox}[Forma de baza]
\begin{lstlisting}[style=codequest]
cat timp (conditie) executa
    instructiuni
sfarsit cat timp
\end{lstlisting}
\end{conceptbox}

\begin{guidebox}[Cheia acestei lectii]
  Pentru lucrul cu cifrele, conditia naturala este \texttt{n != 0}.
\end{guidebox}

\newpage
\SetPageTheme{guided}
\PageHeader{Lectia 5}{Bucla pe cifre}{L5 • P2}

\begin{examplebox}[Schema de baza]
\begin{lstlisting}[style=codequest]
cat timp (n != 0) executa
    uc = n % 10;
    // lucram cu uc
    n = n / 10;
sfarsit cat timp
\end{lstlisting}
\end{examplebox}

\begin{guidebox}[Intrebari ghidate]
  1. Ce reprezinta \texttt{uc} la fiecare pas?

  2. De ce trebuie sa existe instructiunea \texttt{n = n / 10}?

  3. Ce s-ar intampla daca ea ar lipsi?
  \AnswerSpace{5}
\end{guidebox}

\newpage
\SetPageTheme{guided}
\PageHeader{Lectia 5}{Trace pe bucla}{L5 • P3}

\begin{solobox}[Exercitiu semi-ghidat]
  Completeaza tabelul pentru \texttt{n = 5724}:
  \begin{center}
  \begin{tabularx}{\linewidth}{|X|X|X|}
    \hline
    \textbf{Pas} & \textbf{uc} & \textbf{n dupa /10} \\
    \hline
    Start & - & 5724 \\
    \hline
    1 & 4 & 572 \\
    \hline
    2 & & \\
    \hline
    3 & & \\
    \hline
    4 & & \\
    \hline
  \end{tabularx}
  \end{center}
\end{solobox}

\begin{solobox}[Exercitiu independent scurt]
  Fa singur un tabel similar pentru \texttt{n = 9081}.
  \AnswerSpace{4}
\end{solobox}

\newpage
\SetPageTheme{challenge}
\PageHeader{Lectia 5}{Consolidare}{L5 • P4}

\begin{solobox}[Fisa independenta]
  1. Explica in cuvinte cum stie bucla \texttt{cat timp} cand sa se opreasca.

  2. Scrie schema generala prin care prelucram toate cifrele unui numar.

  3. Arata de ce \texttt{n = 0} este un caz special atunci cand vrem sa numaram cifrele.
  \AnswerSpace{8}
\end{solobox}
